\section{State of the art in plasma shape control}
\label{sec:sota_shapectrl}
To define a plasma shape and its position, two key elements are considered: the \emph{centroid} and the \emph{last closed flux surface} (LCFS).

The centroid is the center of the plasma cross-section along the poloidal plane, and can be either \emph{magnetic} or \emph{geometric}; the type of centroid to consider depends on the specific application at hand. The magnetic centroid is usually involved when it is convenient to approximate the plasma as a filament carrying the plasma current; in that case, the magnetic centroid corresponds to this ideal filament revolving around the vessel. The geometric centroid, on the other hand, can be considered when no such approximation is relevant and the entirety of the plasma cross-section within the vessel must be taken into account. When dealing with magnetic control, which is also the case of this work, the magnetic centroid is usually considered.

Poloidal field coils, running along the torus, generate a poloidal field that induces a \emph{poloidal flux} in the vessel. Along the plasma volume, there are points in which this flux has the same value. Examining the plasma from a poloidal cross-section, one can group this points by flux values, thus identifying \emph{isoflux surfaces}. These run outward starting from the plasma centroid, initially as closed surfaces that enclose one another. The last closed one of these surfaces is identified as the LCFS.

Thus, controlling the plasma shape usually involves tracking reference setpoints for the centroid and the LCFS, expressed as a set of control points. The shape control capabilities directly impact operational parameters and plasma-material interactions, allowing for advanced scenarios that maximize performance while preserving machine integrity. This type of control represents a critical tool for achieving optimal plasma conditions while maintaining operational safety margins. To achieve different plasma configurations in terms of shape, the machine must be equipped with both a vessel large enough to contain them and with a set of coils that can generate the necessary magnetic fields.
